% ============================================================================
% templates/report/prueba-de-comprension/prueba-de-comprension.tex: plantilla,
% prueba de comprensión de la audiencia tras una exposición de estudiantes
% ----------------------------------------------------------------------------
% Por qué, instrucciones al expositor y criterio de la rúbrica; después, cada
% una en su página para imprimir, la hoja del expositor y la de la audiencia; al
% final, el conteo del docente, un ejemplo completo y las referencias. SIN diseño:
% lo pone academic-report.
% Uso: scripts/build.sh prueba-de-comprension.tex   (LuaLaTeX)
% ============================================================================
\documentclass{academic-report}
\tipodocumento{Prueba de comprensión}
\universidad{Universidad Nacional de San Cristóbal de Huamanga}
\escuela{Escuela Profesional de Economía}
\curso{Nombre del curso}\docente{Edison Achalma B.Sc. Econ.}\periodo{2026-I}
\titulodocumento{Lo que la audiencia entendió de una exposición}

\begin{document}
\encabezado

\section*{Por qué}
Una exposición se juzga por lo que la audiencia entendió, no por lo que el expositor dijo. Así
se midió en el estudio de Garner y Alley (2013): con 110 estudiantes de ingeniería, las
respuestas escritas de quienes vieron láminas de aserción y evidencia mostraron mejor
comprensión y menos concepciones erróneas que las de quienes vieron láminas con los formatos
por defecto (p.~1564). Además, responder preguntas tras estudiar retiene más que volver a
estudiar cuando se mide a los dos días o a la semana (Roediger y Karpicke, 2006): la prueba
mide al expositor y, de paso, hace aprender a la audiencia.

\noindent\textbf{Fuente:} Garner y Alley (2013); Roediger y Karpicke (2006).

\section*{Instrucciones al expositor}
\begin{enumerate}
  \item Antes de exponer, escribe tu \textbf{mensaje central} en una oración y \textbf{dos o tres
        preguntas} sobre él, con su respuesta esperada y qué se acepta como correcto. Entrégalas
        al docente antes de la sesión; la audiencia no las ve hasta el final.
  \item Cada pregunta pide una idea de tu mensaje, no un detalle de una lámina. Respuesta breve:
        una cifra, una palabra o una oración.
  \item Al terminar, el docente apaga el proyector y reparte la hoja de la audiencia. Se responde
        en tres minutos, sin láminas a la vista y sin apuntes.
  \item La hoja de la audiencia no lleva nombre: mide la exposición, no a quien responde.
  \item El docente cuenta los aciertos y te entrega el porcentaje por pregunta. Una pregunta que
        resultó ambigua se descarta antes de contar, y se dice.
\end{enumerate}

\section*{Cómo entra en la rúbrica de exposiciones}
El promedio de aciertos es el criterio \textbf{Comprensión de la audiencia} de la rúbrica de
exposiciones. Los puntos de cada nivel los fija la rúbrica del curso.

\noindent
\begin{tabularx}{\linewidth}{@{}l c X c@{}}
\toprule
\textbf{Nivel} & \textbf{Aciertos} & \textbf{Qué significa} & \textbf{Pts.} \\\midrule
Destacado  & 75\,\% o más    & La audiencia se lleva el mensaje central y sabe explicarlo. & \\
Logrado    & 50 a 74\,\%     & Se lleva el dato, pero no siempre el porqué.                & \\
En proceso & 25 a 49\,\%     & Se lleva fragmentos; el mensaje central no quedó claro.     & \\
En inicio  & menos de 25\,\% & El mensaje no llegó: revisar estructura y láminas.          & \\
\bottomrule
\end{tabularx}


% ------------------------------------------------------ hoja del expositor ---
\newpage
\section*{Hoja del expositor}
Se entrega al docente antes de la sesión; la audiencia no la ve.

\vspace{0.3em}
\noindent
\begin{tabularx}{\linewidth}{|X|X|X|}
\hline
\textbf{Expositor (código):} & \textbf{Título:} & \textbf{Fecha:} \\[0.5em]
\hline
\multicolumn{3}{|l|}{\textbf{Mensaje central} (una oración):} \\[1.6em]
\hline
\end{tabularx}

\vspace{0.4em}\noindent
\begin{tabularx}{\linewidth}{|c|X|X|X|}
\hline
\textbf{N.º} & \textbf{Pregunta} & \textbf{Respuesta esperada} & \textbf{Se acepta si\ldots} \\
\hline
1 & & & \rule[-4.2cm]{0pt}{4.2cm} \\
\hline
2 & & & \rule[-4.2cm]{0pt}{4.2cm} \\
\hline
3 & & & \rule[-4.2cm]{0pt}{4.2cm} \\
\hline
\end{tabularx}

% ---------------------------------------------------- hoja de la audiencia ---
\newpage
\section*{Hoja de la audiencia}
Sin nombre. Responde sin mirar láminas ni apuntes; si no sabes, déjalo en blanco.

\vspace{0.3em}\noindent
\begin{tabularx}{\linewidth}{|X|X|}
\hline
\textbf{Exposición:} & \textbf{Fecha:} \\[0.5em]
\hline
\end{tabularx}

\vspace{0.5em}\noindent
\begin{tabularx}{\linewidth}{|c|X|}
\hline
\textbf{N.º} & \textbf{Pregunta y tu respuesta} \\
\hline
1 & \rule[-5.4cm]{0pt}{5.4cm} \\
\hline
2 & \rule[-5.4cm]{0pt}{5.4cm} \\
\hline
3 & \rule[-5.4cm]{0pt}{5.4cm} \\
\hline
\end{tabularx}

% ------------------------------------------------------- conteo del docente ---
\newpage
\section*{Conteo del docente}
\noindent
\begin{tabularx}{\linewidth}{|X|X|X|}
\hline
\textbf{Exposición:} & \textbf{Hojas recibidas ($n$):} & \textbf{Fecha:} \\[0.5em]
\hline
\end{tabularx}

\vspace{0.4em}\noindent
\begin{tabularx}{\linewidth}{|c|Y|Y|Y|X|}
\hline
\textbf{N.º} & \textbf{Respuestas} & \textbf{Aciertos} & \textbf{\% de aciertos} & \textbf{Observación} \\
\hline
1 & & & & \\[0.6em]
\hline
2 & & & & \\[0.6em]
\hline
3 & & & & \\[0.6em]
\hline
\multicolumn{3}{|r|}{\textbf{Promedio de las preguntas}} & & \\[0.6em]
\hline
\end{tabularx}
\notatabla{\% de aciertos = aciertos / hojas recibidas × 100; una hoja en blanco en una
pregunta cuenta como respuesta no acertada. El promedio es simple entre las preguntas que no se
descartaron.}

% --------------------------------------------------------------- ejemplo ---
\section*{Ejemplo completo}
Exposición de 5 minutos: \textit{¿Por qué la pobreza de 2024 (27,6\,\%) sigue por encima de la
de 2019 (20,2\,\%)?}

\textbf{Mensaje central.} La pobreza monetaria de 2024 sigue 7,4 puntos por encima de la de 2019
porque el gasto real por persona no volvió a su nivel previo a la pandemia: 903 soles al mes en
2024 frente a 999 en 2019, a precios de 2025 (INEI, 2026, pp.~24, 72). Como pobre es quien gasta
menos que la línea de pobreza, el costo de una canasta básica de alimentos y de otros bienes y
servicios (INEI, 2026, p.~11), menos gasto real significa más personas bajo la línea.

\noindent
\begin{tabularx}{\linewidth}{@{}c >{\raggedright\arraybackslash}X >{\raggedright\arraybackslash}p{0.22\linewidth} >{\raggedright\arraybackslash}p{0.2\linewidth}@{}}
\toprule
\textbf{N.º} & \textbf{Pregunta} & \textbf{Respuesta esperada} & \textbf{Se acepta si\ldots} \\\midrule
1 & ¿Cuántos puntos porcentuales separan la pobreza de 2024 de la de 2019?
  & 7,4 (27,6 menos 20,2). & da entre 7 y 7,5. \\\addlinespace
2 & Según la exposición, ¿qué no ha vuelto a su nivel de 2019 y explica la diferencia?
  & El gasto real por persona. & dice gasto o ingreso real de los hogares. \\\addlinespace
3 & ¿Con qué se compara el gasto de una persona para decir que es pobre?
  & Con la línea de pobreza (el costo de una canasta básica). & nombra la línea o la canasta. \\
\bottomrule
\end{tabularx}

\vspace{0.6em}\noindent
\textbf{Conteo} (datos ilustrativos, 30 hojas recibidas):

\noindent
\begin{tabularx}{\linewidth}{@{}c c c c X@{}}
\toprule
\textbf{N.º} & \textbf{Respuestas} & \textbf{Aciertos} & \textbf{\% de aciertos} & \textbf{Observación} \\\midrule
1 & 30 & 24 & 80\,\% & El dato llegó. \\
2 & 30 & 18 & 60\,\% & Varios responden “la pandemia”, no el gasto. \\
3 & 30 & 12 & 40\,\% & La línea de pobreza se explicó en una lámina recargada. \\\midrule
\multicolumn{3}{@{}r}{\textbf{Promedio de las preguntas}} & \textbf{60\,\%} & Nivel aceptable. \\
\bottomrule
\end{tabularx}
\notatabla{Cifras de pobreza y gasto: INEI (2026, pp.~24, 72). El conteo es un ejemplo, no un
resultado real.}

\textbf{Devolución al expositor.} La audiencia retuvo la cifra, pero no el mecanismo. Para la
próxima, una lámina cuyo titular sea la aserción (“Menos gasto real, más personas bajo la línea”)
y como evidencia el gráfico del gasto de 2019 a 2024.

\section*{Referencias}
{\setlength{\parindent}{0pt}
\hangindent=1.27cm\hangafter=1
Garner, J. K., y Alley, M. (2013). How the design of presentation slides affects audience
comprehension: A case for the assertion-evidence approach. \textit{International Journal of
Engineering Education, 29}(6), 1564–1579.\par
\hangindent=1.27cm\hangafter=1
Instituto Nacional de Estadística e Informática. (2026). \textit{Perú: Evolución de la pobreza
monetaria 2016-2025. Informe técnico}. INEI.\par
\hangindent=1.27cm\hangafter=1
Roediger, H. L., III, y Karpicke, J. D. (2006). Test-enhanced learning: Taking memory tests
improves long-term retention. \textit{Psychological Science, 17}(3), 249–255.
\url{https://doi.org/10.1111/j.1467-9280.2006.01693.x}\par}

\end{document}
